\label{app:speed_training}

In this section we describe the computational overhead of our implementation on a \emph{CIFAR-10} shaped set of random data. In our experimental setup, we use a standard neural network with two convolutional layers followed by max pooling operations and a linear layer. We study how a drop-in replacement of the vanilla PyTorch layers by our chosen Lipschitz constrained layers affects the overall training time of our network. Importantly, our implementation is batch-size independent since it only relies on the values of the weights. Therefore the overall cost of training Lipschitz constrained networks is mitigated when using large batch sizes. 

\begin{table}[h]
  \centering
  \sisetup{
    detect-weight = true,
    detect-family = true,
    table-number-alignment = center
  }
  \begin{tabular}{
    S[table-format=4.0]
    S[table-format=+3.2]
  }
    \toprule
    {Batch size} & {Train time (\%)} \\
    \midrule
    1024 & +20.31 \\
    2048 & +10.26 \\
    \bottomrule
  \end{tabular}
  \caption{Training time on random data of shape \((32,32,3)\) when dropping in the Lipschitz‐constrained layers. For reference, our ResNeXt model of Figure 3 introduces a 9.6\% runtime overhead on TinyImageNet compared to an identical unconstrained model.}
  \label{tab:speed_training}
\end{table}


For more information regarding the computational overhead of Lipschitz constrained architectures w.r.t standard unconstrained architectures, we refer the reader to the comprehensive benchmark of \cite{boissin2025adaptiveorthogonalconvolutionscheme}.

\paragraph{Improvements} In our paper we use a relatively standard implementation of Lipschitz parametrized networks with an orthogonality constraint. Many recent works have focused on providing efficient and expressive Lipschitz network implementations \cite{sll,gloro,cpl} that could be applied in the context of our study. 